\scanpage{137}

\uline{Corollaire} Supposons que l'homomorphisme $M\handwritten{\longrightarrow}M\otimes_{k}M$ induise un homomorphisme $M\otimes_{(A,\handwritten{\Delta}_{A})}(A\otimes_{k}A)\quad M\otimes_{k}M\deleted{\unsure{,|x}}$ \deleted{donc par réduction} qui soit un \uline{isomorphisme} Alors l'homomorphisme $u\handwritten{\mapsto}u.e$ de $A$ dans $M$ est un isomorphisme pour toutes les structures.

Il suffit de prouver que le rang $\inserted{n}$ de $M$ sur $A$ est $1$. Or ce rang de $M\otimes_{k}M$ sur $A\otimes_{k}A$ est $n^{2}$, et on aura $n=n^{2}$, i.e. $n(n-1)=0$ d'où, comme $n\handwritten{\not}=0$, $n=1$. Ceci ne marche que si on sait que le rang est \unsure{fini|mini}. Mais on peut dire aussi que l'application diagonale pour $\overset{\handwritten{-}}{M}$, déduite de celle de $M$ par passage au quotient (NB $\overset{\handwritten{-}}{M}\otimes\overset{\handwritten{-}}{M}=(M\otimes_{k}M)\otimes_{(A\otimes_{k}A)}k$) est \uline{surjective}, donc bijective puisqu'elle est unitaire, donc tout élément de $\overset{\handwritten{-}}{M}\otimes\overset{\handwritten{-}}{M}$ est de la forme $\overset{\handwritten{-}}{e}\otimes x$, ce qui n'est possible que si $\overset{\handwritten{-}}{e}$ engendre $\overset{\handwritten{-}}{M}$.

Ceci montre que la conjecture de locale trivialité \unsure{soulevée|sonlevée} dans mes notes sur les $\otimes$-catégories marche si la bialgèbre $A$ \unsure{envisagëe|envisagée|enfisagëe|enfisagée} est finie sur $k$\inserted{et sa duale est cocomplète,} \uline{pourvu que} l\deleted{a}'algèbre $B$ admette un homomorphisme d'algèbres dans $k$ (ou une extension de $k$), i.e. pourvu que $B$ ne soit pas identique à son idéal des commutateurs.
% archive sheet number: 136
